%
% Teil 3: Verzweigungspunkte und Verzweigungsschnitte am Beispiel
% des halbunendlichen Waermeleiters; Schluesselloch-Kontur (Keyhole)
%
\section{Verzweigungsschnitte und die Schlüsselloch-Kontur}
\label{laplace:subsection:verzweigung}

Wie schnell mehrdeutige Funktionen auftreten, zeigt das Paradebeispiel der
instationären Wärmeleitung.
Ein halbunendlicher Stab $x \geq 0$ mit Anfangstemperatur null wird durch
die normierte Wärmeleitungsgleichung
\begin{equation}
    \frac{\partial^2 U}{\partial x^2}
    =
    \frac{\partial U}{\partial t}
    \label{laplace:eq:waermeleitung}
\end{equation}
beschrieben, wobei am Ende $x = 0$ eine vorgegebene Randtemperatur
$A_0(t)$ eingeprägt wird und die verschwindende Anfangstemperatur formal
durch die Anfangsbedingung
\begin{equation}
    U(x, 0)
    =
    0
    \label{laplace:eq:anfangsbedingung}
\end{equation}
ausgedrückt wird.
Die Laplace-Transformation bezüglich der Zeit macht aus der Zeitableitung
gemäss der Ableitungsregel aus
Tabelle~\ref{laplace:tab:korrespondenzen} eine Multiplikation,
\begin{equation}
    \mathcal{L}\left\{\frac{\partial U}{\partial t}\right\}
    =
    s \, u(x,s) - U(x,0)
    =
    s \, u(x,s),
    \label{laplace:eq:ableitungsregel}
\end{equation}
wobei der Anfangswertterm genau wegen
Gleichung~\eqref{laplace:eq:anfangsbedingung} entfällt.
Aus der partiellen Differentialgleichung wird so eine gewöhnliche
Differentialgleichung für die Bildfunktion
$u(x,s) = \mathcal{L}\{U(x,t)\}$, nämlich
\begin{equation}
    \frac{d^2 u}{d x^2}
    =
    s \, u,
    \label{laplace:eq:bildgleichung}
\end{equation}
mit der allgemeinen Lösung
$u(x,s) = c_1 e^{x\sqrt{s}} + c_2 e^{-x\sqrt{s}}$.
Da die Temperatur für $x \to \infty$ beschränkt bleiben muss, scheidet der
anwachsende Anteil aus; das setzt allerdings $\text{Re}\sqrt{s} > 0$
voraus, denn nur dann wächst $e^{x\sqrt{s}}$ tatsächlich an, während
$e^{-x\sqrt{s}}$ beschränkt bleibt.
Auf dem gleich eingeführten Hauptzweig der Wurzel ist genau das der Fall,
sodass $c_1 = 0$ folgt.
Mit der transformierten Randbedingung
$a_0(s) = \mathcal{L}\{A_0(t)\}$ erhalten wir
\begin{equation}
    u(x,s)
    =
    a_0(s) \, u_0(x,s),
    \qquad
    u_0(x,s)
    =
    e^{-x\sqrt{s}},
    \label{laplace:eq:halbunendlich_loesung}
\end{equation}
wie in \cite{laplace:doetsch1961} im Detail hergeleitet wird.
Die Lösung zerfällt also in die transformierte Randanregung $a_0(s)$ und
eine \emph{Basislösung} $u_0(x,s)$, die am Rand $x = 0$ den Wert eins
annimmt und das reine Ortsprofil beschreibt; dieselbe Aufspaltung wird uns
beim Stab endlicher Länge wieder begegnen.
Die Geometrie des halbunendlichen Gebietes hat uns also zwangsläufig eine
Wurzel $\sqrt{s}$ eingebrockt.

Diese Wurzel ist keine harmlose Funktion, sondern \emph{mehrdeutig}.
Zu jedem $s \neq 0$ existieren zwei Werte $w$ mit $w^2 = s$, und welcher
davon gemeint ist, hängt vom gewählten Zweig ab.
Schreiben wir $s = r e^{i\theta}$ und umlaufen den Ursprung einmal, indem
$\theta$ von $0$ bis $2\pi$ wächst, so geht
$\sqrt{s} = \sqrt{r}\, e^{i\theta/2}$ in
$\sqrt{r}\, e^{i\pi} = -\sqrt{r}$ über:
Die Funktion kehrt nicht zu ihrem Ausgangswert zurück, sondern wechselt
das Vorzeichen \cite{laplace:ablowitz2003}.
Der Ursprung $s = 0$ ist deshalb ein \emph{Verzweigungspunkt} von
$\sqrt{s}$.
Für den Residuensatz ist das fatal, denn er setzt eine eindeutige,
holomorphe Funktion voraus.
Der Ausweg besteht darin, geschlossene Umläufe um den Verzweigungspunkt von
vornherein zu verbieten:
Wir schlitzen die Ebene entlang eines \emph{Verzweigungsschnittes} auf, der
den Verzweigungspunkt $s = 0$ mit dem zweiten Verzweigungspunkt
$s = \infty$ verbindet, und arbeiten auf dem Hauptzweig
$-\pi < \arg s < \pi$.
Der Schnitt liegt dann auf der negativen reellen Achse, und diese Wahl ist
kein Zufall.
Auf dem Hauptzweig gilt nämlich $\text{Re}\sqrt{s} \geq 0$, sodass
$|e^{-x\sqrt{s}}| \leq 1$ beschränkt bleibt, was genau der physikalischen
Beschränktheit der Temperatur entspricht.
Ausserdem bleibt die rechte Halbebene mitsamt dem Bromwich-Pfad
unangetastet.

Nur dürfen wir die D-Kontur aus
Abschnitt~\ref{laplace:subsection:residuensatz} jetzt nicht mehr einfach
zuziehen, denn ihr Kreisbogen würde den Schnitt bei $s = -R$ einmal
durchstossen.
Die Kontur muss dem Schnitt ausweichen:
Sie läuft von oben kommend bis knapp an die negative reelle Achse heran,
folgt ihr oberhalb des Schnittes nach innen, umrundet den
Verzweigungspunkt auf einem kleinen Kreis $\Gamma_\rho$ und kehrt
unterhalb des Schnittes wieder nach aussen zurück.
Das Ergebnis ist die berühmte \emph{Schlüsselloch-Kontur}, die in
Abbildung~\ref{laplace:fig:keyhole_kontur} dargestellt ist.
Innerhalb dieser Kurve ist der Integrand eindeutig und holomorph bis auf
isolierte Pole, sodass der Residuensatz wieder anwendbar ist.

\begin{figure}
    \centering
    \scalebox{0.8}{
\begin{tikzpicture}

% ==============================
% 1. FARBEN DEFINIEREN
% ==============================
\definecolor{plotblue}{RGB}{0, 0, 200}
\definecolor{plotred}{RGB}{220, 30, 30}
\definecolor{plotgray}{RGB}{140, 140, 140}
\definecolor{plotbg}{RGB}{244, 244, 252}
\definecolor{plotorange}{RGB}{230, 140, 0}
\definecolor{plotgreen}{RGB}{0, 150, 80}
\definecolor{plotpurple}{RGB}{140, 60, 180}

% ==============================
% 2. ACHSEN UND SCHNITT
% ==============================
% Verzweigungsschnitt (Wellenlinie)
\draw[thick, plotred, decorate, decoration={snake, amplitude=0.8mm, segment length=3mm}] (-5.1, 0) -- (0, 0);

\draw[thick, -{Stealth[length=3.5mm, width=2.5mm]}] (-5.5, 0) -- (2.5, 0) node[right] {\textbf{Re}$(s)$};
\draw[thick, -{Stealth[length=3.5mm, width=2.5mm]}] (0, -5.0) -- (0, 5.0) node[above] {\textbf{Im}$(s)$};

% Verzweigungspunkt
\fill[plotred] (0,0) circle (2.5pt);

% ==============================
% 3. KONTUR-SEGMENTE
% ==============================
% 1. Bromwich-Pfad (Blau)
\draw[thick, plotblue, decoration={markings, mark=at position 0.6 with {\arrow{Stealth[length=3mm, width=2mm]}}}, postaction={decorate}] 
    (1.2, -4.025) -- (1.2, 4.025);
\draw[thick] (1.2, 0.1) -- (1.2, -0.1) node[below right, inner sep=2pt] {$\boldsymbol{\gamma}$};

% 2. Gamma_R oben (Orange)
\draw[thick, plotorange, decoration={markings, mark=at position 0.5 with {\arrow{Stealth[length=3mm, width=2mm]}}}, postaction={decorate}] 
    (1.2, 4.025) arc (73.39:177.95:4.2);
\node[plotorange, font=\large] at (-3.5, 3.5) {$\Gamma_R$};

% 3. L_+ (Gruen)
\draw[thick, plotgreen, decoration={markings, mark=at position 0.5 with {\arrow{Stealth[length=3mm, width=2mm]}}}, postaction={decorate}] 
    (-4.197, 0.15) -- (-0.37, 0.15);
\node[plotgreen, anchor=south] at (-2.0, 0.65) {$L_+$};

% 4. Gamma_rho (Violett)
\draw[thick, plotpurple, decoration={markings, mark=at position 0.5 with {\arrow{Stealth[length=3mm, width=2mm]}}}, postaction={decorate}] 
    (-0.37, 0.15) arc (157.98:-157.98:0.4);
\node[plotpurple, font=\large] at (0.7, 0.5) {$\Gamma_\rho$};

% 5. L_- (Gruen)
\draw[thick, plotgreen, decoration={markings, mark=at position 0.5 with {\arrow{Stealth[length=3mm, width=2mm]}}}, postaction={decorate}] 
    (-0.37, -0.15) -- (-4.197, -0.15);
\node[plotgreen, anchor=north] at (-2.0, -0.65) {$L_-$};

% 6. Gamma_R unten (Orange)
\draw[thick, plotorange, decoration={markings, mark=at position 0.5 with {\arrow{Stealth[length=3mm, width=2mm]}}}, postaction={decorate}] 
    (-4.197, -0.15) arc (182.05:286.61:4.2);

% ==============================
% 4. ZWEIG-BESCHRIFTUNGEN
% ==============================
\node[plotgray, font=\small, anchor=south] at (-2.0, 0.2) {$s=re^{i\pi},\ \sqrt{s}=i\sqrt{r}$};
\node[plotgray, font=\small, anchor=north] at (-2.0, -0.2) {$s=re^{-i\pi},\ \sqrt{s}=-i\sqrt{r}$};

% ==============================
% 5. POLE (Kreuze)
% ==============================
\newcommand{\pole}[3][plotred]{
    \draw[thick, #1, line cap=round] (#2-0.15,#3-0.15) -- (#2+0.15,#3+0.15);
    \draw[thick, #1, line cap=round] (#2-0.15,#3+0.15) -- (#2+0.15,#3-0.15);
}

% Pole auf imag. Achse
\pole{0.0}{2.2}
\node[plotred, right] at (0.15, 2.2) {$i\omega$};

\pole{0.0}{-2.2}
\node[plotred, right] at (0.15, -2.2) {$-i\omega$};

% ==============================
% 6. LEGENDE
% ==============================
\node[draw, thick, rounded corners=3pt, fill=white, inner xsep=5pt, inner ysep=4pt, anchor=west] at (3.8, 0) {
    \renewcommand{\arraystretch}{1.2}
    \begin{tabular}{@{}c l@{}}
        
        \tikz[baseline=-0.6ex]{
            \fill[plotred] (0,0) circle (2pt);
        } & Verzweigungspunkt \\
        
        \tikz[baseline=-0.6ex]{
            \draw[thick, plotred] (-0.45, 0) .. controls (-0.35, 0.2) and (-0.25, -0.2) .. (-0.15, 0) .. controls (-0.05, 0.2) and (0.05, -0.2) .. (0.15, 0) .. controls (0.25, 0.2) and (0.35, -0.2) .. (0.45, 0);
        } & Verzweigungsschnitt \\
        
        \tikz[baseline=-0.6ex]{
            \draw[thick, plotred, line cap=round] (-0.1,-0.1) -- (0.1,0.1) (-0.1,0.1) -- (0.1,-0.1);
        } & Polstelle \\
        
        \tikz[baseline=-0.6ex]{
            \draw[thick, plotblue, ->-=0.55] (-0.3,0) -- (0.3,0);
        } & Bromwich-Pfad \\
        
        \tikz[baseline=-0.6ex]{
            \draw[thick, plotorange, ->-=0.55] (-0.3,0) -- (0.3,0);
        } & Kreisbogen $\Gamma_R$ \\
        
        \tikz[baseline=-0.6ex]{
            \draw[thick, plotgreen, ->-=0.55] (-0.3,0) -- (0.3,0);
        } & Ufer $L_+, L_-$ \\
        
        \tikz[baseline=-0.6ex]{
            \draw[thick, plotpurple, ->-=0.55] (-0.3,0) -- (0.3,0);
        } & Kreisbogen $\Gamma_\rho$ \\
        
    \end{tabular}
};

\end{tikzpicture}
}

    \caption{Die Schlüsselloch-Kontur für Bildfunktionen mit
    Verzweigungspunkt bei $s = 0$.
    Der Verzweigungsschnitt (rote Wellenlinie) auf der negativen reellen
    Achse wird von den beiden horizontalen Wegen $L_+$ und $L_-$
    eingerahmt, der kleine Kreis $\Gamma_\rho$ umgeht den
    Verzweigungspunkt.
    Innerhalb der Kontur ist $\sqrt{s}$ eindeutig, und der Residuensatz
    erfasst nur die eingeschlossenen Pole, hier $s = \pm i\omega$.}
    \label{laplace:fig:keyhole_kontur}
\end{figure}

Wie diese Kontur konkret arbeitet, zeigt das klassische Beispiel aus
\cite{laplace:doetsch1961}, bei dem das Stabende mit der Temperatur
$A_0(t) = \cos \omega t$ angeregt wird.
Wegen $\mathcal{L}\{\cos \omega t\} = \frac{s}{s^2 + \omega^2}$ lautet die
Bildfunktion
\begin{equation}
    u(x,s)
    =
    \frac{s}{s^2 + \omega^2} \, e^{-x\sqrt{s}},
    \label{laplace:eq:keyhole_beispiel}
\end{equation}
mit einfachen Polen bei $s = \pm i\omega$ und dem Verzweigungspunkt bei
$s = 0$.
Der Residuensatz auf der Schlüsselloch-Kontur liefert $2\pi i$ mal die
Residuen an den beiden eingeschlossenen Polen.
Im Grenzübergang verschwinden die grossen Bögen nach dem Lemma von Jordan,
denn $\frac{s}{s^2+\omega^2}$ fällt wie $\frac{1}{s}$ ab und
$|e^{-x\sqrt{s}}| \leq 1$ stört dabei nicht.
Auch der kleine Kreis $\Gamma_\rho$ trägt nichts bei, weil der Integrand
dort beschränkt ist und die Weglänge gegen null schrumpft.
Die Residuen selbst rechnen wir in zwei Schritten aus.
Am einfachen Pol $s = i\omega$ liefert die Grenzwertformel zunächst
\begin{equation}
    \operatorname{Res}\bigl(u(x,s)e^{st}, i\omega\bigr)
    =
    \lim_{s \to i\omega}
    (s - i\omega) \,
    \frac{s \, e^{st - x\sqrt{s}}}{(s - i\omega)(s + i\omega)}
    =
    \frac{i\omega \, e^{i\omega t - x\sqrt{i\omega}}}{2 i\omega}
    =
    \frac{1}{2} \, e^{i\omega t - x\sqrt{i\omega}}.
    \label{laplace:eq:keyhole_grenzwert}
\end{equation}
Die verbleibende Wurzel werten wir in Polarform aus.
Wegen $i = e^{i\pi/2}$ ist
$\sqrt{i\omega}
= \sqrt{\omega} \, e^{i\pi/4}
= \sqrt{\omega} \, \frac{1 + i}{\sqrt{2}}
= \sqrt{\tfrac{\omega}{2}} \, (1 + i)$,
und am Pol $s = -i\omega$ ergibt sich völlig analog der konjugierte
Wert.
Zusammengefasst lauten die beiden Residuen also
\begin{equation}
    \operatorname{Res}\bigl(u(x,s)e^{st}, \pm i\omega\bigr)
    =
    \frac{1}{2} \,
    e^{-x\sqrt{\omega/2}} \,
    e^{\pm i\left(\omega t - x\sqrt{\omega/2}\right)},
    \label{laplace:eq:keyhole_residuen}
\end{equation}
und ihre Summe ergibt den reellen Ausdruck
$e^{-x\sqrt{\omega/2}} \cos\bigl(\omega t - x\sqrt{\omega/2}\bigr)$.

Der eigentliche Clou passiert aber am Schnitt.
Oberhalb gilt $s = r e^{i\pi}$ und damit $\sqrt{s} = i\sqrt{r}$, unterhalb
dagegen $s = r e^{-i\pi}$ und $\sqrt{s} = -i\sqrt{r}$.
Die Integranden auf $L_+$ und $L_-$ sind also verschieden, die beiden
Wegintegrale heben sich \emph{nicht} auf, und genau ihre Differenz bleibt
als reelles Integral übrig.
Sortiert man alle Beiträge, so entsteht die Lösung
\begin{equation}
    U(x,t)
    =
    e^{-x\sqrt{\omega/2}}
    \cos\Bigl(\omega t - x\sqrt{\tfrac{\omega}{2}}\Bigr)
    -
    \frac{1}{\pi}
    \int_0^\infty
    e^{-rt} \, \frac{r}{r^2 + \omega^2} \, \sin\bigl(x \sqrt{r}\bigr)
    \, dr.
    \label{laplace:eq:keyhole_loesung}
\end{equation}
Diese Formel erzählt die ganze Physik des Problems.
Der erste Term stammt von den Polen und beschreibt den eingeschwungenen
Zustand:
Eine Temperaturwelle dringt in den Stab ein, wird dabei exponentiell
gedämpft und läuft der Anregung mit wachsender Tiefe zunehmend hinterher.
Der zweite Term stammt vom Verzweigungsschnitt und ist ein transienter
Einschwingvorgang, der wegen des Faktors $e^{-rt}$ für grosse Zeiten
abklingt.
Während isolierte Pole diskrete Beiträge $e^{s_k t}$ liefern, wirkt der
Schnitt wie ein \emph{Kontinuum} von Abklingraten $e^{-rt}$, über das
integriert werden muss.

Damit drängt sich eine Frage auf.
Der Verzweigungsschnitt kam über die Wurzel $\sqrt{s}$ ins Spiel, und die
Wurzel kam über das halbunendliche Gebiet.
Was passiert, wenn der Stab eine endliche Länge $l$ besitzt?
Die Antwort ist verblüffend und der schönste Aha-Moment dieser Arbeit.
